% !TeX root = ../../../research-log.tex

\subsection*{YD $\leftrightarrow$ double}

For a finite dimensional Hopf algebra $H$, it is known that $\mathcal{YD}^H_H \cong \mathrm{Mod}_{D(H)}$.
We can deduce the cross relation on $D(H)$ form the Yetter--Drinfel'd condition.

\begin{definition}
    Let $V$ be a right $H$-module and $H$-comodule simultaneously.
    Then, $V$ is csalled a Yetter--Drinfel'd module, shortly YD-module, if 
    \begin{align} \label{eqn:YD}
        v_0 \triangleleft h_1 \otimes v_1h_2 = (v \triangleright h_2)_0 \otimes h_1 (v \triangleleft h_2)_1.
    \end{align}
\end{definition}

For a right $H$-module $V$, let define a left $H^*$ action by
\[
    f \triangleright v := \langle f, v_1 \rangle v_0.
\]

For any $f \in H^*$, by applying $\langle f, - \rangle$ to the second factor of \eqref{eqn:YD}, we have
\begin{align*}
v_0 \triangleleft h_1 \langle f, v_1h_2 \rangle & = v_0 \triangleleft h_1 \langle f_1, v_1 \rangle \langle f_2, h_2 \rangle \\
& = (f_1 \triangleright v) \triangleleft h_1 \langle f_2, h_2 \rangle
\end{align*}
and
\begin{align*}
(v \triangleright h_2)_0 \langle f, h_1 (v \triangleleft h_2)_1 \rangle & = (v \triangleright h_2)_0 \langle f_1, h_1 \rangle \langle f_2, (v \triangleleft h_2)_1 \rangle \\
& = \langle f_1, h_1 \rangle (f_2 \triangleright (v \triangleleft h_2)).
\end{align*}
Therefore, the YD condition \eqref{eqn:YD} is equivalent to
\[
(f_1 \triangleright v) \triangleleft h_1 \langle f_2, h_2 \rangle = \langle f_1, h_1 \rangle (f_2 \triangleright (v \triangleleft h_2))
\]
for all $v \in V$, and hence
\[
f_1 h_1 \langle f_2, h_2 \rangle = \langle f_1, h_1 \rangle h_2 f_2.
\]

From the above arguments, we may have the following.

\begin{proposition}
    \[
    \mathcal{YD}^H_H \cong D(H) \mathrm{-Mod}_{H\text{-rat}}
    \]
\end{proposition}

We verify the following proposition in \cite{Kashaev21}.

\begin{proposition}
    $D(H)^\circ$ has a dual quasitriangular structure with the $R$-form $\mathscr{R}: D(H)^\circ \otimes D(H)^\circ \rightarrow k$ defined by
    \[
    \mathscr{R}(x \otimes y) := \langle x, (\jmath \circ \imath^\circ)(y) \rangle
    \]
\end{proposition}
\begin{proof}
    
\end{proof}